\section{Related Work}
\label{sec:related_work}

\noindent\textbf{RL Training Frameworks.}~Recent systems have significantly improved the scalability of reinforcement learning for foundation models. A line of works~\cite{verl,siiRL,areal,streamRL,miroRL,roll} provide system support for distributed RL training, including scalable rollout-update pipelines, asynchronous execution, and integration with model-parallel training backends. More recently, several frameworks~\cite{longrl,verl,slime,easyr1} have begun to support RL post-training for vision-language models. For example, LongRL~\cite{longrl} targets RL for long-video workloads, EasyR1~\cite{easyr1} extend general RLHF infrastructure to multimodal models. These systems make VLM RL training increasingly practical, but they do not explicitly address the workload imbalance introduced by multimodal inputs. %In particular, image-text and video-text samples can vary substantially in visual token counts, sequence lengths, and preprocessing costs, leading to pronounced imbalance across devices and pipeline stages during rollout and training.

% \noindent\textbf{RL Training Frameworks.}~Reinforcement learning has become a central paradigm for aligning and enhancing large language models (LLMs). Recent frameworks such as veRL~\citep{verl}, siiRL~\citep{siiRL}, AReal~\citep{areal}, StreamRL~\citep{streamRL}, MiroRL~\citep{miroRL}, and ROLL~\citep{roll} provide system-level support for distributed RL training. These frameworks focus on issues such as asynchronous rollout-update decoupling, scalable data pipelines, and integration with model-parallel training backends. LongRL is proposed to support long-context VLM RL training, but it does not address the load imbalance challenges when sequence lengths vary drastically~\cite{longrl_issue3,longrl}. 

\noindent\textbf{VLM Training Frameworks.}~The emergence of multimodal LLMs has motivated training frameworks such as DistTrain~\citep{DistTrain}, DistMM~\citep{DistMM}, LongVILA~\citep{LongVILA}, and VeOmni~\citep{veomni}. These systems propose optimizations for heterogeneous architectures combining vision towers and language backbones. Despite the advances, existing VLM frameworks primarily target pretraining or supervised fine-tuning. They do not explicitly address the unique challenges of RL training.

% HybridDP method. FlexSP, HotSpa, Hydrallus, ByteScale, WLB-LLM, LongVILA.   KimiVL, InternVL3.5. KeyeVL, GLM4.5V.
\noindent\textbf{Load Balancing for Large Model Training.}~A line of work focuses on load balancing techniques for efficient large model training. Classic methods rely on sequence bucketing and packing, which sort samples by length and allocate them across GPUs to minimize padding~\citep{KimiVL,KeyeVL,GLM-4_5V,InternVL3_5}. Recent works such as FlexSP~\citep{FlexSP}, HotSPa~\citep{HotSPa}, Hydraulis~\citep{Hydraulis}, ByteScale~\citep{ByteScale}, and WLB-LLM~\citep{WLB-LLM} mitigate load imbalance via heterogeneous parallelism and dynamic reconfiguration across DP ranks or gradient steps. Zeppelin~\citep{Zeppelin} addresses training imbalance by jointly optimizing attention-specific sequence partitioning and layout remapping between attention and linear layers.

\section{Conclusion}
\label{sec_discussion}

In this work, we present \SysName, a holistic system that addresses the unique system-level challenges of reinforcement learning for large Vision-Language Models. By systematically analyzing the bottlenecks across the RL pipeline, we identify critical inefficiencies in both data loading and workload balancing that hinder scalability and hardware utilization. \SysName~introduces ShadowLoader to eliminate I/O and memory bottlenecks on the single controller, and FlexUlysses to achieve fine-grained, shard-level load balancing across GPUs. Our efficient scheduling algorithm and dynamic execution engine further maximize overlap between computation and communication. Combining both components achieves high throughput even under extreme data heterogeneity.

% \textcolor{blue}{
% \noindent\textbf{Comparison with Dynamic Sequence Packing (DSP).}
% Dynamic sequence packing, aka bucketing, addresses sequence-level padding but fails to mitigate stragglers when a single sequence is an outlier (e.g., a 32K video context). In such cases, the DP rank holding the long sequence becomes the system bottleneck. FlexUlysses operates at sub-sequence granularity, sharding outliers across multiple GPUs to balance both compute and memory, avoiding the OOM and throughput regression of sequence-level methods.
% }

% \textcolor{blue}{
% \noindent\textbf{Applicability to Non-Attention Architectures.}
% While FlexRL is evaluated on attention-based VLMs, its core mechanism—partitioning along the head dimension—is applicable to non-attention or hybrid architectures (e.g., SSMs, Linear Attention) that retain a multi-head structure (e.g., DeepSeek R1, Nemotron, Kimi). The scheduling heuristic can naturally incorporate the cost model of any multi-head operator.
% }

% \textcolor{blue}{
% \noindent\textbf{Limitations.}
% FlexRL introduces extra All-to-All communication, which may become dominant as the sharding degree increases or in clusters with limited inter-node bandwidth. Furthermore, while FlexRL is broadly applicable to multi-head architectures, its gains are most pronounced in regimes with extreme sequence length heterogeneity.
% }
